\documentclass[10pt]{article}

\usepackage[utf8]{inputenc}

\usepackage[T1]{fontenc}

\usepackage{amsmath}

\usepackage{amsfonts}

\usepackage{amssymb}

\usepackage[version=4]{mhchem}

\usepackage{stmaryrd}

\usepackage{graphicx}

\usepackage[export]{adjustbox}

\graphicspath{ {./images/} }



\begin{document}

Теорема Виноградова - Бомбьери о распределении простых чисел в арифметич. прогрессиях в среднем также приводит к решению Т. п. При этом предположение о справедливости расширенной гипотезы Римана заменяется фактически теоремами типа большого решета. Лит.: [1] Л и н н и к Ю. В., Дисперсионный метод в бинарных аддитивных задачах, [Л.], 1961 ; [2] Б р е д их и н Б. М., "Успехи матем. наук», 1965, т. 20 , в. 2 , с. 89-130; [.] П р а х а р к., Распределение простых чисел, пер. с нем.,

М., 1967.



ТИХОНОВА ТЕОРЕМА о бикомпактности произведения: топологич. произведение любого множества бикомпактных прастрансте бикомпактно. Это одна из основных теорем общей топологии; установлена А. Н. Тихоновым в 1929. Она играет весьма существенную и часто ключевую роль в построении практически всех разделов общей топологии и во многих ее применениях . В частности, Т. т. имеет основное значение для построения бикомпактных расширений вполне регулярных $T_{1}$-пространств (т. е. т И р а н с т в). С ее помощью строится расширение Стоуна - Чеха произвольного тихоновского пространства. T. т. позволяет указать стандартные бикомпактные пространства - обобщенные канторовы дисконтинуумы $D \tau$, являющиеся произведениями дискретных двоеточий в количестве $\tau$ и тихоновские кубы $I^{\tau}-$ произведения $\tau$ экземпляров обычного отрезка $I$ числовой прямой. В качестве $\tau$ здесь может фигурировать любой кардинал. Значение обобщенных канторовых дисконтинуумов $D \tau$ и тихоновских кубов $I^{\tau}$ связано прежде всего с тем, что они являются универсальными объектами: каждый нульмерный бикомпакт гомеоморфен замкнутому подпространству нек-рого $D^{\tau}$ и каждый бикомпакт гомеоморфен замкнутому подпространству нек-рого $I \tau$.



Т. т. применяется при доказательстве непустоты предела обратного спектра из бикомпактных пространств, при построении теории абсолютов, в теории бикомпактных групі. Если же иметь в виду опосредованные ее применения, то почти вся общая топология попадает в сферу действия этой Т. т. Так же трудно перечислить прямые и опосредованные применения Т. т. в других областях математики. Практически они встречаются всюду, где важную роль играет понятие компактности, - в частности в функциональном анализе (банаховы пространства в слабой топологии, меры на топологич. пространствах), в общей теории оптимального управления И т. Д.



Лит.: [1] К е л ли Дж., Общая топология, пер. с англ., 2 изд., М., 1981; [2] А р хан г е ль с ки й А. В., П о н ом а p e в В. И., Основы общей топологии в задачах и упражне-



$\begin{array}{rr}\text { ТИХОНОВСКИИ } & \text { КУБ }- \text { топологич. пронгельский. } \\ \text { ТИведение }\end{array}$ $\tau$ экземпляров обычного отрезка $I$ действительной прямой, где $\tau-$ произвольный кардинал; обозначается I $\tau$. Т. к. введен А. Н. Тихоновым в 1929. Если $\tau=n-$



\includegraphics[max width=\textwidth, center]{2023_07_09_641d80896d98b2ccd8abg-1}

$n$-мерном евклидовом пространстве, топология к-рого порождена метрикой скалярного произведения. Если



\includegraphics[max width=\textwidth, center]{2023_07_09_641d80896d98b2ccd8abg-1(1)}

меоморфен аильбертову кирпичу. При $\tau_{1} \neq \tau_{2}$. Т. к. $I^{\tau_{1}}$ и $I^{\tau_{2}}$ не гомеоморфны между собой: если $\tau-$ бесконечный кардинал, то $\tau$ есть вес пространства $I^{\tau}$, а если $\tau=n$ - натуральное число, то $n$ - размерность пространства $I^{n}$. Два свойства Т. к. I $\tau$ особенно важны: бикомпактность каждого из них, независимо от $\tau$, и их универсальность по отношению ко вполне регулярным $T_{1}$-пространствам веса не большего, чем $\tau$ : каждое такое пространство гомеоморфно нек-рому подпространству пространства $I \tau$. Бикомпактные хаусдорфовы пространства, вес к-рых не превосходит $\tau$, гомеоморфны замкнутым подпространствам тихоновского куба $I \tau$. Таким образом, всего двух операций - операции топо- логич. произведения и операции перехода к замкнутому подпространству - достаточно для того, чтобы получить из одного стандартного и весьма простого топологич. пространства - отрезка - любой бикомпакт. Примечательным следствием бикомпактности Т. к. является бикомпактность единичного шара в сопряженном к банахову пространству, наделенном слабой топологией сопряженного. Универсальность Т. к. и простота определения делает их важными стандартными объектами общей топологии. Однако топологич. строение Т. к. далеко не тривиально. В частности, куб $I^{\mathcal{C}}$, где $\mathfrak{c}$ - мощность континуума, сепарабелен, хотя состоит из $2^{\mathfrak{C}}$ точек; вес его равен с. Неожиданный факт: число Суслина каждого Т. к. I $\tau$ счетно, независимо от $\tau$, т. е. каждое семейство попарно непересекающихся открытых множеств в $I \tau$ счетно. Хотя в Т. к. есть много сходящихся последовательностей, этих последних не хватает для того, чтобы описать прямо оператор замыкания в Т. к.



ТИХОНОВСКОЕ ПРОИЗВЕДЕНИЕ С е м е Й с т В а то по ло ги ч е с к и х п р о с т р а н с т в- то же что топологическое произведение семейства топологич. пространств. Понятие Т. П. введено А. Н. Тихоновым (1929).



ТИХОНОВСКОЕ ПРОСТРАНСТВО - топологическое пространство, в к-ром каждое конечное множество замкнуто и для всякого замкнутого множества $\boldsymbol{P}$ и любой не принадлежащей $P$ точки $x$ найдется непрерывная вещественная функция $f$ на всем пространстве, к-рая принимает значение 0 в точке $x$ и значение 1 во всех точках множества $P$. Класс Т. п. совпадает с классом вполне регулярных $T_{1}$-пространств. В Т. П. любые две различные точки отделимы непересекающимися окрестностями - т. е. выполняется аксиома отделимости Хаусдорфа, но не всякое Т. п. нормально. А. Н. Тихонов (1929) охарактеризовал Т. п., как подпространства бикомпактных хаусдорфовых пространств.



Лит.: [1] А л е к с а н д р о в П. С., Введение в теорию множеств и общую топологию, М., 1977; [2] А р х а н г е л ьс к и й А. В., П о н о м а p е в В. И., Основы общей тополо-



ТКАНЕЙ ГЕОМЕТРИЯ - раздел дифференциальной геометрии, в к-ром изучаются нек-рые семейства линий и поверхностей - т. н. ткани (плоские, пространственные, многомерные).



П л о с к о й $p-$ т к а н ь ю наз. область плоскости, в к-рой заданы $p$ (обычно $p \geqslant 3$ ) семейств достаточно гладких линий со свойствами: 1) через каждую точку области проходит точно по одной линии каждого семейства; 2) линии разных семейств имеют не более одной общей точки. Пример: три семейства прямых, параллельных сторонам равностороннего треугольника, образуют 3-ткань (р е г у л я р н у ю, или п р ав и л ь н у ю т к а н ь).



Основным предметом изучения в Т. г. являются свойства, инвариантные при диффференциально топологич. преобразованиях. Ткани наз. эквивалентными, если они (локально или глобально) диффеоморфны. При $p=2$ ткань диффеоморфна ткани, образованной двумя семействами параллельных прямых (такие ткани наз. сетями). При $p=3$ ткань уже в общем случае ведиффеоморфна ни трем семействам параллельных прямых (т. е. не является пा е с т и у г о л ь в о й т к а н ь ю), ни трем семействам прямых вообще (т. е. не является с п р я м л я е мо й т к а н ь ю). Условие шестиугольности ткани в геометрич. форме состоит в выполнении замыкания условия. У словие спрямляемости ткани не может быть записано в обозримом виде; его исследованием занимаются в связи с проблемами номографии.



П рост т ан с т в ен ны е к ри волинейны $\mathbf{~ т ~ к ~ а ~ и ~ о б р а з у ю т с я ~} p$ семсйствами кривых в





\end{document}